\documentclass[11pt]{article}
\usepackage[utf8]{inputenc}
\usepackage[letterpaper,top=0cm, margin=0.85in]{geometry}

\usepackage{textcmds} %more symbols
\usepackage{fontspec} %more fonts


%for math
\usepackage{amsmath, amssymb, amsfonts} %standard
\usepackage{youngtab} % makes squares for math diagrams
\usepackage{microtype} %% <-- added
%-----------------------------------------------------------           

%\usepackage{sectsty}
%for lists and numbers
\usepackage{enumitem}
%-----------------------------------------------------------

% Doc setting
\usepackage[english]{babel} % Replace `english' with e.g. `spanish' to change the document language
\usepackage{setspace} %to set spacing bw words and lines
\usepackage{changepage}
% \setlength\parindent{0pt}

%footer
\usepackage{fancyhdr}
\usepackage{lastpage}

% \fancyhf{0pt} % sets both header and footer to nothing
% \renewcommand{\footrulewidth}{0pt}
% \renewcommand{\headrulewidth}{0pt} %remove headerline

% \fancyfoot[RE,RO]{\thepage}
\fancyfoot[C]{MATH 201 | \emph{Shubhro Gupta}}
\renewcommand{\headrulewidth}{0pt}% Default \headrulewidth is 0.4pt
\renewcommand{\footrulewidth}{0.4pt}% Default \footrulewidth is 0pt
\fancyfoot[L]{Assignment 1}
\fancyfoot[R]{\thepage}
\pagestyle{fancy}
%-----------------------------------------------------------

%for pictures and graphs
\usepackage{graphicx} %add image
\usepackage{adjustbox}

\usepackage{pgfplots} %for graphing plotting
\pgfplotsset{compat=1.18, width=10cm}
%-----------------------------------------------------------

%for code
\usepackage{verbatim}
\usepackage{listings}
\usepackage{fancyvrb} %for coding blocks
%\usepackage{algorithm}
%\usepackage{algpseudocode} %for pseudocode
%\usepackage{algorithm, algpseudocode}

%\usepackage{lstfiracode} %firacode
\usepackage[framemethod=tikz]{mdframed} %adding background to lstlisting
\usepackage[ruled,vlined,boxed]{algorithm2e} %for pseudocode lines



%for colors and links
\usepackage[colorlinks = true,
            linkcolor = blue,
            urlcolor  = blue,
            citecolor = blue,
            anchorcolor = blue]{hyperref}
\usepackage[many]{tcolorbox}  % for colored boxes
\usepackage{color} % to get colors
\usepackage{xcolor} %more colors options and flexibility
\usepackage{transparent}


%-----------------------------------------------------------------------------
%custom commands

%code
\newcommand{\problem
}[2]{
\begin{mdframed}
    \textsc{question} \textbf{#1} \hfill #2
\end{mdframed}
}
\newcommand{\codecap}[2]{{\vspace{4pt}{\emph{#1}}} \hfill \href{#2}{Link to the code\ }\vspace{25pt}}
\newcommand{\code}[1]{{\texttt{#1}}}

%math
\newcommand{\bigo}[1]{$O(#1)$ }
\newcommand{\thetan}[1]{$\theta(#1)$}
% \newcommand{\vector}[1]{$\overrightarrow{#1}$}}o
\newcommand{\contradiction}{%
\begin{tikzpicture}[rotate=45,x=0.5ex,y=0.5ex]
\draw[line width=.2ex] (0,2) -- (3,2) (0,1) -- (3,1) (1,3) -- (1,0) (2,3) -- (2,0);
\end{tikzpicture}
}

\newcommand{\vecset}[2]{\{ {#1}_1, {#1}_2, {#1}_3,  \dots,  {#1}_{#2}\}}

%display
\newcommand{\link}[3][blue]{\href{#2}{\color{#1}{#3}}}%
\newcommand{\inlink}[1]{\underline{\emph{\link[black]{#1}{#1}}}}


%header
\newcommand{\heading}[5]{
\hrule ~\\~\\
\begin{large}
\noindent\emph{#1}\smallskip ~\\
Professor #3 \hfill Week #2 \smallskip ~\\
\textbf{Shubhro Gupta} \hfill Due #4 ~\\
\end{large} \medskip ~\\
{\emph{Collaborators: #5}}~\\
\hrule
\vspace{50pt}
~\\
}

% \newcommand\dunderline[3][-1pt]{{%
%   \sbox0{#3}%
%   \ooalign{\copy0\cr\rule[\dimexpr#1-#2\relax]{\wd0}{#2}}}}


%-----------------------------------------------------------------------------
%title
\usepackage{algpseudocode}
\begin{document}

\heading{Real Analysis}{1}{Rishi Vyas}{Tuesday, July 30, 2024}{none}

\problem{A}{4 Points}
\textbf{Definition.} \emph{Let $x \in \mathbb{R}$. If $n \in \mathbb{N}$, we define $x^n:=\overbrace{x \ldots x}^{n \text { times }}$. By convention, for $n=0$ we interpret this as defining $x^0:=1$.}~\\
\\
\textbf{To Prove.}  Let $x, y, z \in \mathbb{R}$.

\begin{enumerate}
	\item If 0' is an element of $\mathbb{R}$ such that $0'+x=x$ for all $x \in \mathbb{R}$, then $0'=0$.
	      \begin{align*}
		      0' + x        & = x                                             \\
		      0' + x + (-x) & = x + (-x)                                      \\
		      0' + 0        & = 0        \qquad \text{(Additive Inverse (4))} \\
		      0'            & = 0                                             \\
		      \square
	      \end{align*}
	\item If 1' is an element of $\mathbb{R}$ such that $1' \cdot x=x$ for all $x \in \mathbb{R}$, then $1'=1$.
	      \begin{align*}
		      1' \cdot x              & = x                                                   \\
		      1' \cdot x \cdot x^{-1} & = x \cdot x^{-1}                                      \\
		      1' \cdot 1              & = 1        \qquad \text{(Multiplicative Inverse (8))} \\
		      1'                      & = 1                                                   \\
		      \square
	      \end{align*}
	\item $x \neq 0$ and $x y=x z$. Prove that $y=z$. Deduce that if $x y=1$ then $y=x^{-1}$.
	      \begin{align*}
		      x y                  & = x z                  \\
		      ((x)^{-1} \cdot x) y & = ((x)^{-1} \cdot x) z \\
		      y                    & = z                    \\
		      \square
	      \end{align*}
	      \begin{align*}
		      x y                  & = 1                \\
		      ((x)^{-1} \cdot x) y & = (x)^{-1} \cdot 1 \\
		      \square
	      \end{align*}
	\item $-0=0$
	      \begin{align*}
		      0           & = 0           \\
		      0 \times -1 & = 0 \times -1 \\
		      0 = - 0                     \\
	      \end{align*}
	\item If $x \neq 0$, then $x^{-1} \neq 0$ and $\left(x^{-1}\right)^{-1}=x$.
	      \begin{align*}
		      \text{Suppose }x^{-1} & = 0                                                    \\
		      x \cdot x^{-1}        & = x \cdot 0                                            \\
		      1                     & = 0         \qquad \text{(Multiplicative inverse (8))}
	      \end{align*}
	      $1 \neq 0$ so $x^{-1} \neq 0$.
	\item $(-x) \times(-y)=x y$.
	      \begin{align*}
		      (-x) \times(-y)  & = (-x)(-y) + 0y                                          \\
		                       & = (-x)(-y) + (x + (-x))y                                 \\
		                       & = (-x)(-y) + xy + (-x)y \qquad \text{(Distributive Law)} \\
		                       & = ((-x)(-y) + (-x)y) + xy                                \\
		                       & = (-x)(-y + y) + xy                                      \\
		                       & = (-x)(0) + xy                                           \\
		                       & = 0 + xy                                                 \\
		      (-x) \times (-y) & = xy
	      \end{align*}
	\item If $x \neq 0$ and $n \in \mathbb{N}$, then $(-x)^{-1}=-\left(x^{-1}\right)$ and $\left(x^{-1}\right)^n=\left(x^n\right)^{-1}$.
	\item If $x \neq 0$ and $y \neq 0$, then $x y \neq 0$.
	      \begin{align*}
		      \text{Suppose } x y    & = 0             \\
		      x \cdot x^{-1} \cdot y & = 0\cdot x^{-1} \\
		      1 \cdot y              & = 0\cdot x^{-1} \\
		      y                      & = 0
		      \\
		      \contradiction
	      \end{align*}

\end{enumerate}

\problem{B}{3 Points}
\textbf{Definition.} Let $x \in \mathbb{R}$, $x \neq 0.$ If $n \in \mathbb{N}$,
~\\
\noindent \textbf{To Prove.} Let $x, y, z \in \mathbb{R}$.

\begin{enumerate}
	\item If $x<y$ and $z>0$, then $x z<y z$.
	\item If $x<0$ then $x^{-1}<0$.
	\item If $x, y \geq 0$ and $n \in \mathbb{Z}_{>0}$, prove that $x \leq y$ if and only if $x^n \leq y^n$. Deduce that $x<y$ if and only if $x^n<y^n$.

	      \problem{D}{4 Points}
	      \begin{enumerate}
		      \item Let $a, b \in \mathbb{R}$.
		            \begin{enumerate}
			            \item Consider the set $\{a, b\}$. Prove that the supremum of the set $\{a, b\}$ is equal to $a$ if $a \geq b$, and is equal to $b$ if $b \geq a$.
			            \item Prove that
			                  $$
				                  \sup (\{a, b\})=\frac{a+b+|a-b|}{2} .
			                  $$
			            \item  Formulate and prove a variant of the above formula describing the infimum of the set $\{a, b\}$.

		            \end{enumerate}

		      \item Let $a_1, \ldots, a_n$ be elements of $\mathbb{R}$. Prove that
		            $$ \sup \left(\left\{a_1, \ldots, a_n\right\}\right)=\sup \left(\left\{\sup \left(\left\{a_1, \ldots, a_{n-1}\right\}\right), a_n\right\}\right) . $$
		            Deduce that $\sup \left(\left\{a_1, \ldots, a_n\right\}\right) \in\left\{a_1, \ldots, a_n\right\}$.
	      \end{enumerate}



	      \textbf{Definition.} \emph{If $a_1, \ldots, a_n$ are elements of $\mathbb{R}$, define the maximum of these elements by $\max \left\{a_1, \ldots, a_n\right\}:=\sup \left(\left\{a_1, \ldots, a_n\right\}\right)$. From the previous question, we then observe that $\max \left\{a_1, \ldots, a_n\right\} \in\left\{a_1, \ldots, a_n\right\}$. We can define the minimum of $a$ collection of elements in $\mathbb{R}$ in an analogous manner using inf instead of sup.}
	      \\

	      \textbf{Definition.} \emph{Let $S \subseteq \mathbb{R}$. We say that $S$ is bounded if it is bounded above and bounded below.}

	      \problem{E}{3 Points}
	      Let $S$ be a subset of $\mathbb{R}$.
	      \begin{enumerate}
		      \item     Show that $S$ is bounded if and only if there exists an $m \in \mathbb{R}$ such that $|x| \leq m$ for all $x \in S$.
		      \item  Deduce that the following three statements are equivalent:
		            \begin{enumerate}
			            \item The $S$ is bounded.
			            \item There exists $m \in \mathbb{R}$ such that $S \subseteq[-m, m]$.
			            \item There exist $a, b \in \mathbb{R}$ with $a \leq b$ such that $S \subseteq[a, b]$.
		            \end{enumerate}
	      \end{enumerate}

	      Let $S \subseteq \mathbb{R}$.
	      (1) Let $x, y \in \mathbb{R}$, and suppose $x$ and $y$ both satisfy the definition of an infimum of $S$. Show that $x=y$.
	      (2) Let $x \in \mathbb{R}$. Show that $x=\inf (S)$ if and only if it satisfies the following two conditions:
	      a.) The element $x$ is a lower bound for $S$.
	      b.) Given $\epsilon>0$, there exists some $y \in S$ such that $y<x+\epsilon$.
	      (3) Let $x \in \mathbb{R}$. Show that $x=\inf (S)$ if and only if it satisfies the following two conditions:
	      a.) The element $x$ is a lower bound for $S$.
	      b.) Given $n \in \mathbb{Z}_{>0}$, there exists some $y \in S$ such that $y<x+\frac{1}{n}$.
\end{enumerate}

\end{document}
